\documentclass[11pt,a4paper]{article}
\usepackage[utf8]{inputenc}
\usepackage[T1]{fontenc}
\usepackage[spanish]{babel}
\usepackage{geometry}
\usepackage{booktabs}
\usepackage{array}
\usepackage{longtable}
\usepackage{hyperref}
\usepackage{enumitem}
\usepackage{xcolor}
\usepackage{fancyhdr}
\usepackage{titlesec}
\geometry{margin=2.2cm}
\setlength{\headheight}{14pt}
\hypersetup{colorlinks=true,linkcolor=blue,urlcolor=blue}
\pagestyle{fancy}
\fancyhf{}
\lhead{Lab-02 Part-02}
\rhead{Reproduccion Experimental}
\cfoot{\thepage}
\setlist[itemize]{nosep,leftmargin=1.5em}
\setlist[enumerate]{nosep,leftmargin=1.7em}
\titleformat{\section}{\large\bfseries}{}{0em}{}
\titleformat{\subsection}{\normalsize\bfseries}{}{0em}{}

\title{Informe Tecnico\\Parte II: Reproduccion Experimental del Paper}
\author{CIB12 Aprendizaje de Maquina y Mineria de Datos}
\date{7 de abril de 2026}

\begin{document}
\maketitle

\section{Introduccion}
Este informe documenta la reproduccion experimental del paper \textit{Bluetooth low energy indoor positioning: A fingerprinting neural network approach}. El objetivo fue implementar el pipeline principal del articulo con herramientas equivalentes a las reportadas por los autores, ejecutar el modelo M2 y contrastar los resultados obtenidos con los valores publicados.

La reproduccion se organizo dentro de \texttt{02-laboratorios/Lab-02/Part-02}, separando claramente codigo fuente, datos, resultados y evidencias. El entregable incluye codigo reproducible, documentacion tecnica, metricas de ejecucion y un analisis comparativo con el paper.

En terminos sencillos, el problema que resuelve el paper es el siguiente: dada una lectura de intensidad de senal BLE captada por varios receptores, estimar en que posicion del espacio interior se encuentra una persona o un dispositivo. Esto es util porque en interiores el GPS no suele funcionar bien, por lo que se necesitan tecnicas alternativas de localizacion.

\section{Paper Reproducido}
\begin{itemize}
\item Titulo: \textit{Bluetooth low energy indoor positioning: A fingerprinting neural network approach}
\item Autores: Alberto Ferrero-Lopez, Antonio Javier Gallego, Miguel Angel Lozano
\item Revista: \textit{Internet of Things}
\item Ano: 2025
\item DOI: \texttt{10.1016/j.iot.2025.101565}
\end{itemize}

El trabajo propone un sistema de posicionamiento indoor con senales BLE RSSI y \textit{fingerprinting}. El modelo central de esta reproduccion fue M2, una red neuronal para regresion de coordenadas $(x,y)$, evaluada con distancia euclidiana sobre trayectorias reales.

Conviene aclarar tres terminos del paper:
\begin{itemize}
\item \textbf{BLE}: es Bluetooth Low Energy, una tecnologia inalambrica de bajo consumo.
\item \textbf{RSSI}: es la intensidad con la que llega la senal. No da la distancia exacta, pero si ofrece una pista indirecta sobre la cercania entre emisor y receptor.
\item \textbf{Fingerprinting}: en vez de calcular la posicion con geometria pura, se aprende un patron. Cada posicion del entorno queda representada por una ``huella'' de intensidades RSSI observadas por varios sensores.
\end{itemize}

\section{Metodologia Reproducida}
\subsection{Verificacion metodologica del paper}
Se contrastaron las secciones metodologicas del articulo con la implementacion realizada. Los elementos verificados fueron:
\begin{itemize}
\item Separacion entre fase \textit{offline} y fase \textit{online}.
\item Construccion de \textit{fingerprints} RSSI a partir del dataset BLE publico.
\item Uso de dos variantes del dataset: \texttt{with\_empty\_values} y \texttt{without\_empty\_values}.
\item Entrenamiento del modelo M2 con TensorFlow.
\item Evaluacion final mediante distancia euclidiana en metros.
\end{itemize}

La verificacion metodologica fue importante porque la reproduccion no consiste solo en ``entrenar cualquier red''. El objetivo fue respetar la misma logica experimental del paper. Por eso se reviso que la entrada del modelo, el preprocesamiento, la forma de dividir el problema y la metrica final fueran coherentes con el articulo original.

\subsection{Recursos utilizados}
\begin{itemize}
\item Lenguaje: Python 3.11.9
\item Framework principal: TensorFlow 2.15.1
\item Librerias de apoyo: NumPy 1.26.4, Pandas 2.2.3, SciPy 1.11.4, scikit-learn 1.4.2
\item Dataset: \texttt{Position-Annotated-BLE-RSSI-Dataset}
\item Repositorios base: \texttt{models\_fingerprint\_positioning} y \texttt{rssi\_capturing\_filtering\_library}
\item Repositorio del laboratorio: \href{https://github.com/eduardo202020/cib12-aprendizaje-maquina-mineria-datos}{GitHub del curso}
\end{itemize}

\subsection{Parametros reproducidos}
Se respetaron los principales parametros descritos en el paper:
\begin{itemize}
\item \texttt{Sreads\_min = 2}
\item \texttt{Svalid\_min = 10} para \texttt{with\_empty\_values}
\item \texttt{Svalid\_min = 12} para \texttt{without\_empty\_values}
\item Filtro \texttt{Max}
\item Ventana temporal entre \texttt{0.5 s} y \texttt{2.5 s}
\item Valor \texttt{100} para representar lecturas faltantes
\item Perdida \texttt{MSE}, salida lineal, \texttt{batch\_size = 256}, \texttt{early stopping} y optimizador Adam
\end{itemize}

Estos parametros tienen un sentido practico:
\begin{itemize}
\item \texttt{Sreads\_min} y \texttt{Svalid\_min} controlan cuantas lecturas minimas se aceptan antes de construir un \textit{fingerprint}.
\item El filtro \texttt{Max} intenta quedarse con la lectura mas representativa o mas fuerte dentro de una ventana temporal.
\item El valor \texttt{100} marca ausencia de lectura valida. En otras palabras, no significa una senal real, sino ``dato faltante''.
\item La perdida \texttt{MSE} mide que tan lejos queda la prediccion de las coordenadas reales. Mientras menor sea, mejor esta ajustando el modelo durante el entrenamiento.
\item \texttt{Early stopping} detiene el entrenamiento cuando seguir entrenando ya no mejora el error de validacion, evitando sobreajuste.
\end{itemize}

\subsection{Implementacion}
El pipeline implementado quedo distribuido en los siguientes scripts:
\begin{itemize}
\item \texttt{src/prepare\_datasets.py}: genera los \textit{fingerprints} offline y online.
\item \texttt{src/train\_m2.py}: entrena el modelo M2 y guarda el modelo y el escalador.
\item \texttt{src/evaluate\_m2.py}: evalua el modelo con trayectorias online y calcula distancia euclidiana.
\item \texttt{src/run\_pipeline.py}: orquesta el flujo completo.
\end{itemize}

La idea general del pipeline es: primero transformar datos crudos en ejemplos utiles, luego entrenar el modelo con esos ejemplos y finalmente medir si el modelo generaliza cuando se enfrenta a trayectorias que no vio antes.

\section{Proceso de Reproduccion}
La reproduccion siguio esta secuencia:
\begin{enumerate}
\item Clonado de repositorios auxiliares y descarga del dataset.
\item Creacion del entorno compatible \texttt{.venv311}.
\item Generacion de CSV procesados para ambas variantes del problema.
\item Entrenamiento de M2 con \texttt{with\_empty\_values}.
\item Entrenamiento de M2 con \texttt{without\_empty\_values}.
\item Evaluacion sobre trayectorias reales de la fase online.
\item Barrido exploratorio de semillas para revisar la sensibilidad de \texttt{with\_empty\_values}.
\end{enumerate}

Cada etapa cumple una funcion distinta:
\begin{itemize}
\item En la fase \textit{offline} se aprende. Aqui el sistema ya conoce la posicion real de cada lectura y usa esa informacion para entrenar.
\item En la fase \textit{online} se pone a prueba. Aqui el modelo recibe nuevas lecturas y debe estimar la posicion sin ayuda.
\item El barrido de semillas no cambia la metodologia del paper, pero permite observar si pequeñas diferencias en la inicializacion del entrenamiento afectan el resultado final.
\end{itemize}

Como evidencia de la preparacion del dataset se obtuvieron:
\begin{itemize}
\item \texttt{213558} filas offline y \texttt{662} filas online para \texttt{with\_empty\_values}.
\item \texttt{141713} filas offline y \texttt{486} filas online para \texttt{without\_empty\_values}.
\end{itemize}

La diferencia entre ambas variantes aparece porque \texttt{without\_empty\_values} elimina ejemplos incompletos, mientras que \texttt{with\_empty\_values} conserva mas muestras pero introduciendo el marcador de dato faltante. Por eso una variante tiene mas datos y la otra tiene datos mas ``limpios''.

\section{Resultados}
\subsection{Metricas obtenidas}
\begin{center}
\begin{tabular}{>{\raggedright\arraybackslash}p{4.2cm}cccc}
\toprule
Escenario & Media (m) & Mediana (m) & Min (m) & Max (m) \\
\midrule
M2 with\_empty\_values & 2.3690 & 2.1413 & 0.0743 & 15.9757 \\
M2 with\_empty\_values seed21 & 2.3169 & 1.9804 & 0.1214 & 16.1974 \\
M2 without\_empty\_values & 2.0750 & 1.8430 & 0.0517 & 8.1161 \\
\bottomrule
\end{tabular}
\end{center}

La corrida \texttt{seed21} se incluyo como evidencia de sensibilidad del entrenamiento a la inicializacion. No reemplaza la corrida base, pero permite observar una mejora controlada en el escenario con valores vacios.

Las metricas principales son la media y la mediana del error en metros:
\begin{itemize}
\item \textbf{Media}: promedio del error total. Es sensible a casos muy malos.
\item \textbf{Mediana}: valor central del error. Suele representar mejor el comportamiento tipico del modelo.
\item \textbf{Minimo y maximo}: muestran el mejor y el peor caso observados.
\end{itemize}

Por ejemplo, en \texttt{without\_empty\_values} la media fue \texttt{2.0750 m}. Eso significa que, en promedio, la posicion predicha quedo a poco mas de dos metros de la posicion real. La mediana fue \texttt{1.8430 m}, lo que indica que la mitad de los errores estuvo por debajo de esa distancia.

\subsection{Resumen de entrenamiento}
\begin{center}
\begin{tabular}{>{\raggedright\arraybackslash}p{4.2cm}ccc}
\toprule
Escenario & Epocas & Loss final & Val\_loss final \\
\midrule
with\_empty\_values & 60 & 0.01735 & 0.03107 \\
with\_empty\_values seed21 & 55 & 0.01737 & 0.02942 \\
without\_empty\_values & 30 & 0.00032 & 0.01743 \\
\bottomrule
\end{tabular}
\end{center}

En esta tabla, \texttt{loss} es el error sobre los datos usados para entrenar y \texttt{val\_loss} es el error sobre un subconjunto de validacion que el modelo no usa para ajustarse directamente. La \texttt{val\_loss} es especialmente importante porque indica si el modelo esta aprendiendo patrones utiles o si solo esta memorizando el conjunto de entrenamiento.

\section{Analisis Comparativo}
\subsection{Comparacion con el paper}
\begin{center}
\begin{tabular}{>{\raggedright\arraybackslash}p{4.2cm}cccc}
\toprule
Escenario & Paper media & Reprod. media & Paper mediana & Reprod. mediana \\
\midrule
with\_empty\_values & 2.0368 & 2.3690 & 1.8199 & 2.1413 \\
with\_empty\_values seed21 & 2.0368 & 2.3169 & 1.8199 & 1.9804 \\
without\_empty\_values & 2.0559 & 2.0750 & 1.7745 & 1.8430 \\
\bottomrule
\end{tabular}
\end{center}

La comparacion debe interpretarse asi: mientras mas cerca esten las columnas de ``Reprod.'' respecto a las del paper, mas fiel fue la reproduccion. No se espera necesariamente igualdad exacta, porque intervienen semillas aleatorias, detalles de implementacion y posibles diferencias no descritas por completo en el articulo.

\subsection{Reflexion}
La reproduccion de \texttt{without\_empty\_values} fue muy cercana al paper: la diferencia en media fue de solo \texttt{0.0191 m} y en mediana de \texttt{0.0685 m}. Esto respalda que el pipeline implementado es metodologicamente consistente con el articulo.

En cambio, \texttt{with\_empty\_values} quedo por encima de los resultados publicados. La mejor corrida exploratoria redujo la brecha, pero no logro igualar el paper. Las causas mas probables son:
\begin{itemize}
\item diferencias en el preprocesamiento exacto de los \textit{fingerprints},
\item sensibilidad a la inicializacion aleatoria del entrenamiento,
\item diferencias no completamente especificadas en artefactos intermedios,
\item y posibles variaciones menores entre librerias o versiones del entorno.
\end{itemize}

Desde el punto de vista conceptual, esto sugiere que el escenario con valores faltantes es mas inestable. Cuando el modelo recibe muchas entradas ``vacías'' o marcadas con el valor 100, la relacion entre patron RSSI y posicion se vuelve mas dificil de aprender. En cambio, cuando el dataset elimina esos casos, la red recibe ejemplos mas consistentes y la prediccion mejora.

\section{Documentacion y Evidencias}
Esta seccion muestra evidencia visible dentro del propio informe. No se limita a enumerar nombres de archivos: incluye extractos del contenido generado durante la ejecucion.

\begin{center}
\small
\begin{tabular}{>{\raggedright\arraybackslash}p{3.9cm}p{9.0cm}}
\toprule
Archivo o carpeta & Evidencia que aporta \\
\midrule
\texttt{src/prepare\_datasets.py} & Script que genera los \textit{fingerprints} offline y online. \\
\texttt{src/train\_m2.py} & Script que entrena M2 y genera \texttt{model.keras}, \texttt{scaler.pkl} y \texttt{training\_summary.json}. \\
\texttt{src/evaluate\_m2.py} & Script que calcula la distancia euclidiana y produce los JSON de evaluacion. \\
\texttt{training\_summary.json} & Resume epocas ejecutadas, \texttt{loss} y \texttt{val\_loss} del entrenamiento. \\
\texttt{m2\_...evaluation.json} & Resume cantidad de muestras y estadisticas de error en metros. \\
\texttt{metricas-reproduccion.csv} & Tabla compacta con las metricas finales incorporada a la entrega. \\
\texttt{comandos-ejecutados.txt} & Traza resumida de los comandos principales del pipeline. \\
\bottomrule
\end{tabular}
\end{center}

\subsection*{Extracto del archivo de comandos}
El archivo \texttt{evidencias/base/comandos-ejecutados.txt} conserva la traza de ejecucion. A continuacion se muestra un extracto real incluido en la entrega:

\begin{quote}
\ttfamily\footnotesize
powershell -ExecutionPolicy Bypass -File 01-organizacion-cib12\textbackslash 02-laboratorios\textbackslash Lab-02\textbackslash Part-02\textbackslash scripts\textbackslash bootstrap\_env.ps1\\
python 01-organizacion-cib12\textbackslash 02-laboratorios\textbackslash Lab-02\textbackslash Part-02\textbackslash src\textbackslash prepare\_datasets.py --fingerprint-name lab02\_ble\_ferrero\\
python 01-organizacion-cib12\textbackslash 02-laboratorios\textbackslash Lab-02\textbackslash Part-02\textbackslash src\textbackslash train\_m2.py --fingerprint-name lab02\_ble\_ferrero --variant with\_empty\_values\\
python 01-organizacion-cib12\textbackslash 02-laboratorios\textbackslash Lab-02\textbackslash Part-02\textbackslash src\textbackslash evaluate\_m2.py --fingerprint-name lab02\_ble\_ferrero --variant without\_empty\_values
\end{quote}

Este bloque permite demostrar que la reproduccion no fue solo teorica. Aqui se observa la secuencia real del trabajo: primero preparar el entorno, luego preparar el dataset, despues entrenar el modelo y finalmente evaluarlo.

\subsection*{Extracto del archivo CSV de metricas}
El archivo \texttt{evidencias/base/metricas-reproduccion.csv} resume las metricas finales. Su contenido principal es:

\begin{center}
\begin{tabular}{>{\raggedright\arraybackslash}p{4.1cm}cccc}
\toprule
Escenario & Media & Mediana & Epocas & Val\_loss \\
\midrule
with\_empty\_values & 2.3690 & 2.1413 & 60 & 0.03107 \\
with\_empty\_values seed21 & 2.3169 & 1.9804 & 55 & 0.02942 \\
without\_empty\_values & 2.0750 & 1.8430 & 30 & 0.01743 \\
\bottomrule
\end{tabular}
\end{center}

Este CSV es util porque concentra en un solo lugar los resultados numericos mas importantes. Sirve tanto para documentar el experimento como para comparar rapidamente variantes del mismo modelo.

\subsection*{Extracto del JSON de evaluacion}
Tambien se generaron archivos JSON en \texttt{resultados/comparativas/metricas}. El siguiente bloque reproduce el contenido esencial de \texttt{without\_empty\_values\_evaluation.json}, que corresponde al escenario mas cercano a los resultados del paper:

\begin{quote}
\ttfamily\footnotesize
\{\\
"fingerprint\_name": "lab02\_ble\_ferrero",\\
"variant": "without\_empty\_values",\\
"distance\_stats": \{\\
\hspace*{1em}"count": 486,\\
\hspace*{1em}"min": 0.0517,\\
\hspace*{1em}"median": 1.8430,\\
\hspace*{1em}"mean": 2.0750,\\
\hspace*{1em}"max": 8.1161\\
\}\\
\}
\end{quote}

Aqui se observa que el archivo JSON no guarda solo un numero final, sino varias estadisticas del error. Eso es importante porque una sola metrica puede ocultar comportamientos extremos. Por ejemplo, la media resume el desempeno general, pero el maximo muestra cual fue el peor caso observado.

\subsection*{Extracto del resumen de entrenamiento}
El archivo \texttt{training\_summary.json} del escenario \texttt{without\_empty\_values} muestra que el entrenamiento se ejecuto y convergio con los siguientes valores:

\begin{quote}
\ttfamily\footnotesize
\{\\
"fingerprint\_name": "lab02\_ble\_ferrero",\\
"variant": "without\_empty\_values",\\
"epochs\_ran": 30,\\
"final\_loss": 0.00031697,\\
"final\_val\_loss": 0.01742656\\
\}
\end{quote}

Este archivo ayuda a explicar que el modelo efectivamente fue entrenado y no solo evaluado. El numero de epocas indica cuantas veces el algoritmo recorrio los datos de entrenamiento, mientras que \texttt{final\_loss} y \texttt{final\_val\_loss} resumen el nivel de error alcanzado al final del proceso.

\subsection*{Archivos efectivamente generados}
Durante la ejecucion se obtuvieron, entre otros, los siguientes archivos de salida:
\begin{itemize}
\item \texttt{resultados/experimentos/M2/base/with\_empty\_values/model.keras}
\item \texttt{resultados/experimentos/M2/base/without\_empty\_values/model.keras}
\item \texttt{resultados/experimentos/M2/base/without\_empty\_values/training\_summary.json}
\item \texttt{resultados/comparativas/metricas/M2/base/with\_empty\_values\_evaluation.json}
\item \texttt{resultados/comparativas/metricas/M2/base/without\_empty\_values\_evaluation.json}
\item \texttt{evidencias/base/metricas-reproduccion.csv}
\item \texttt{evidencias/base/comandos-ejecutados.txt}
\end{itemize}

Los artefactos grandes de entrenamiento y los datos procesados permanecen fuera de Git mediante \texttt{.gitignore}, pero fueron usados en la ejecucion real y quedan representados en este informe por sus rutas, sus metricas y los archivos de evidencia versionados dentro de \texttt{Part-02}.

\section{Conclusiones}
Se cumplio la reproduccion experimental del paper implementando el modelo M2 con Python y TensorFlow, siguiendo la logica metodologica reportada por los autores. La documentacion del proceso, los resultados y el analisis comparativo quedaron integrados dentro de \texttt{Part-02}.

La conclusion principal es que la reproduccion fue exitosa en el escenario \texttt{without\_empty\_values}, donde los resultados quedaron muy cercanos a los del paper. El escenario \texttt{with\_empty\_values} presento una brecha mayor, aunque la exploracion de semillas mostro que parte de esa diferencia puede deberse a la sensibilidad del entrenamiento.

En otras palabras, el laboratorio permitio comprobar que la propuesta del paper es reproducible y que su desempeno no depende solo de una afirmacion teorica. Al mismo tiempo, tambien mostro una leccion importante de aprendizaje de maquina: dos experimentos aparentemente iguales pueden diferir cuando cambian detalles pequenos del preprocesamiento, de los datos faltantes o de la inicializacion del entrenamiento. Esa observacion forma parte del aprendizaje tecnico de esta reproduccion.

\end{document}
